\documentclass[11pt]{article}
\usepackage[margin=1in]{geometry}
\usepackage{amsmath,amssymb,amsthm}
\usepackage{xurl}
\usepackage{hyperref}
\sloppy
\let\oldus\_ \renewcommand{\_}{\oldus\allowbreak}
\newtheorem{theorem}{Theorem}
\newtheorem{lemma}[theorem]{Lemma}
\theoremstyle{definition}
\newtheorem{definition}[theorem]{Definition}
\theoremstyle{remark}
\newtheorem{remark}[theorem]{Remark}
\newcommand{\HKd}{\mathrm{HKd}}

\title{Conjecture 00000003888 is false:\\ the Hilbert--Kunz density function of $(k[x,y,z],(x,y,z))$\\ is $x^2/2$ on $[0,1)$, hence not piecewise linear}
\author{Jackmeson1}
\date{}

\begin{document}
\maketitle

\section{The conjecture and the reading}

\textbf{Statement (English).} \emph{Definition: The Hilbert--Kunz density function is the limit of relative eHK as a
function of parameter x. Conjecture: The HK density function of a fixed ideal is piecewise linear with finitely many
breakpoints; conversely, there exist explicit one-parameter ideal families whose breakpoint counts grow unboundedly.}
The Chinese text says the same: the HK density function of a fixed ideal is piecewise linear with finitely many
breakpoints, and conversely there is an explicit one-parameter family of ideals whose breakpoint counts grow without
bound.

\textbf{Reading.} The conjecture is a conjunction (Clause 1) $\wedge$ (Clause 2). We refute Clause 1, which refutes the
conjunction whatever Clause 2 means. We do not address Clause 2. The text does not restrict the ring or its dimension,
so Clause 1 says: for every pair $(R,I)$ for which the HK density function is defined, $\HKd(R,I)$ is piecewise linear
with finitely many breakpoints. We exhibit one pair, $R=k[x,y,z]$ with the standard grading and $I=(x,y,z)$, over any
field $k$ of characteristic $p>0$, for which this fails. For curves (standard graded rings of dimension $2$), piecewise
linearity is a theorem of Trivedi (quoted below). Our counterexample has dimension $3$, which the text does not exclude.

\section{Definitions}

\textbf{Source.} V.~Trivedi, \emph{Hilbert--Kunz density function and Hilbert--Kunz multiplicity},
\url{https://arxiv.org/abs/1510.03294}. The quotes below were retrieved from the arXiv abstract and from the ar5iv
HTML of the paper (\url{https://ar5iv.labs.arxiv.org/html/1510.03294}).
\begin{itemize}
\item Abstract: ``For a pair $(M, I)$, where $M$ is finitely generated graded module over a standard graded ring $R$
of dimension $d$, and $I$ is a graded ideal with $\ell(R/I) < \infty$, we introduce a new invariant $HKd(M, I)$ called
the Hilbert-Kunz density function.''
\item Introduction: ``$HKd(M,I)(x)=f(x)=\lim_{n\to\infty}g_{n}(x)=\lim_{n\to\infty}f_{n}(x),$ where
$f_{n}(x)=(1/q^{d-1})\ell(M/I^{[q]}M)_{\lfloor xq\rfloor}$''. Notations 2.1: ``Fix $n\in\mathbb{N}$ and denote
$q=p^{n}$''.
\item Frobenius power: ``$I^{[p^{n}]}=n$-th Frobenius power of $I$ = the ideal generated by $p^{n}$-th powers of
elements of $I$''.
\item Dimension 2: ``the density function $HKd(R,{\bf m})$, is a piecewise linear polynomial with rational
coefficients''.
\end{itemize}

\begin{definition}[Lean: \texttt{frobeniusPower}, \texttt{gradedPiece}, \texttt{hkApprox}, \texttt{hkDensity}]
Let $k$ be a field of characteristic $p>0$, $R=k[x_\sigma]$ with $|\sigma|=d$ and the standard grading, and $I$ an ideal.
$I^{[q]}$ is the ideal generated by $\{f^q: f\in I\}$. The degree-$m$ piece $(R/J)_m$ is the image of the degree-$m$
homogeneous polynomials in $R/J$, and $\ell((R/J)_m)=\dim_k (R/J)_m$; it is $0$ for $m<0$. Then
\[ f_n(x)=\frac{1}{q^{d-1}}\,\dim_k\bigl(R/I^{[q]}\bigr)_{\lfloor xq\rfloor},\qquad q=p^n,\qquad
\HKd(R,I)(x)=\lim_{n\to\infty} f_n(x). \]
In Lean, $p=\texttt{ringChar}\ k$, $d=\texttt{Nat.card}\ \sigma$ (and \texttt{ringKrullDim\_eq} proves
$\dim k[x_\sigma]=|\sigma|$), and $\HKd$ is \texttt{limUnder atTop}, which equals the limit wherever it exists. This
is Trivedi's definition for $M=R$. A module of length $\ell$ over $R$ concentrated in one degree is a
$k=R_0$-vector space of dimension $\ell$.
\end{definition}

\begin{definition}[Lean: \texttt{IsPiecewiseLinearOn}, \texttt{IsPiecewiseLinearPartition}]
$f:\mathbb{R}\to\mathbb{R}$ is \emph{piecewise linear with finitely many breakpoints on $s$} if there is a finite set
$B$ such that $f$ is affine ($f(x)=\mu x+c$) on every open interval contained in $s\setminus B$. Values at breakpoints
are unconstrained and continuity is not required. The partition form says there are
$a=t_0<t_1<\dots<t_k=b$ with $f$ affine on every $[t_i,t_{i+1}]$. Lemma
\texttt{IsPiecewiseLinearPartition.isPiecewiseLinearOn} proves that the partition form on $[a,b]$ implies
\texttt{IsPiecewiseLinearOn} on $[a,b]$, with $B=\{t_0,\dots,t_k\}$. ``Linear'' pieces are read as affine, which is the
weaker condition. So the refutation also covers the reading with pieces of the form $\mu x$.
\end{definition}

\section{The proof}

Fix a field $k$ of characteristic $p>0$, $R=k[x,y,z]$ ($d=3$) and $\mathfrak m=(x,y,z)$
(Lean: \texttt{maxIdeal}).

\begin{lemma}[\texttt{maxIdeal\_isHomogeneous}, \texttt{finrank\_quotient\_maxIdeal}]
$\mathfrak m$ is homogeneous, and $\dim_k R/\mathfrak m=1$. So $(R,\mathfrak m)$ is in the scope of the definition.
\end{lemma}
\begin{proof}
It is generated by the homogeneous elements $x,y,z$. The kernel of $f\mapsto f(0)$ is $\mathfrak m$: a polynomial
lies in the monomial ideal $\mathfrak m$ iff every monomial in its support is divisible by some variable, that is,
iff its constant coefficient is $0$. So $R/\mathfrak m\cong k$.
\end{proof}

\begin{lemma}[\texttt{frobeniusPower\_maxIdeal}]
For $q=p^n$, $\mathfrak m^{[q]}=(x^q,y^q,z^q)$.
\end{lemma}
\begin{proof}
$\supseteq$ is clear. For $\subseteq$, induct on $f\in\mathfrak m=\mathrm{span}(x,y,z)$. The generators satisfy
$x_i^q\in(x^q,y^q,z^q)$, and $0^q=0$. In characteristic $p$, $(a+b)^q=a^q+b^q$ (\texttt{add\_pow\_char\_pow}), and
$(ra)^q=r^qa^q$.
\end{proof}

\begin{lemma}[\texttt{hilbertFunction\_monomialIdeal}, \texttt{hilbertFunction\_Gq}]
For $0\le m<q$, $\dim_k (R/(x^q,y^q,z^q))_m=\binom{m+2}{2}$.
\end{lemma}
\begin{proof}
For a monomial ideal, $\dim_k(R/J)_m$ is the number of degree-$m$ monomials not divisible by any generator. This is
proved in Lean via \texttt{mem\_ideal\_span\_monomial\_image}, reusing the accepted solution of conjecture 00000000539.
If $m<q$, no degree-$m$ monomial is divisible by $x_i^q$. So all $\binom{m+2}{2}$ degree-$m$ monomials are standard
(stars and bars, \texttt{card\_finsuppAntidiag\_nat\_eq\_choose}).
\end{proof}

\begin{theorem}[\texttt{tendsto\_hkApprox}, \texttt{hkDensity\_maxIdeal}]
For $0\le x<1$, $f_n(x)\to x^2/2$ as $n\to\infty$. Hence $\HKd(R,\mathfrak m)(x)=x^2/2$ on $[0,1)$.
\end{theorem}
\begin{proof}
Let $q=p^n$ and $m=\lfloor xq\rfloor$. Then $0\le m\le xq<q$. By the lemmas,
$f_n(x)=\binom{m+2}{2}/q^2=\tfrac12(m/q+2/q)(m/q+1/q)$. Since $x-1/q<m/q\le x$ and $q\to\infty$, we get
$m/q\to x$ and $1/q\to0$.
\end{proof}
This agrees with Trivedi's Example 3.1 (retrieved), stated for $R=k[X_0,\dots,X_d]$ and its graded maximal ideal:
``$HKd({\mathbb{P}}^{d}_{k})(x)=x^{d}/d!$, for $0\leq x<1$''. Here $d=2$. Our Lean proof does not use this example.

\begin{theorem}[\texttt{not\_affine}, \texttt{not\_isPiecewiseLinearOn}]
If $f=x^2/2$ on $(0,1)$, then $f$ is not affine on any open interval $(a,b)\subseteq(0,1)$ with $a<b$. Hence $f$ is not
piecewise linear with finitely many breakpoints on any $s\supseteq(0,1)$.
\end{theorem}
\begin{proof}
Let $u=(3a+b)/4$, $v=(a+b)/2$ and $w=(a+3b)/4$. An affine $f$ has $f(u)+f(w)-2f(v)=0$, but here this quantity is
$(b-a)^2/16>0$. Given a finite $B$, pick $x\in(0,1)\setminus B$. The set $(0,1)\setminus B$ is open, so it contains
an interval $(x-\varepsilon,x+\varepsilon)$, on which $f$ would have to be affine.
\end{proof}

\begin{theorem}[\texttt{not\_clause1}, \texttt{conjecture\_false}]
Let $k$ be a field of characteristic $p>0$. Clause 1, restricted to the polynomial rings $k[x_1,\dots,x_n]$ and their
homogeneous ideals of finite colength, is false: $\HKd(k[x,y,z],\mathfrak m)$ is not piecewise linear on $\mathbb{R}$.
Hence $\neg(\text{Clause 1}\wedge Q)$ for every proposition $Q$.
\end{theorem}
Further variants are also proved. \texttt{hkDensity\_not\_isPiecewiseLinearOn} works on any domain $s\supseteq(0,1)$,
such as $[0,\infty)$ or $[0,3]$. \texttt{hkDensity\_not\_isPiecewiseLinearPartition} covers the partition form on any
$[a,b]\supseteq[0,1]$. \texttt{limit\_not\_isPiecewiseLinearOn} applies to every function that is the pointwise limit
of $f_n$ on $(0,1)$, so it does not depend on the \texttt{limUnder} convention.

\section{Lean formalization}
File \texttt{lean/Conjecture3888/Basic.lean} (Lean 4, Mathlib v4.33.1, only \texttt{import Mathlib}). Main theorem:
\begin{verbatim}
theorem conjecture_false (K : Type*) [Field K] (p : N) [Fact p.Prime]
    [CharP K p] (Clause2 : Prop) : not (Clause1 K /\ Clause2)
-- Clause1 K : for all n and I : Ideal (MvPolynomial (Fin n) K),
--   I homogeneous -> Module.Finite K (R / I) ->
--   IsPiecewiseLinearOn (hkDensity I) Set.univ
\end{verbatim}
\textbf{Axiom audit.} \texttt{\#print axioms} reports only \texttt{propext}, \texttt{Classical.choice} and
\texttt{Quot.sound} for \texttt{conjecture\_false}, \texttt{not\_clause1},
\texttt{hkDensity\_not\_isPiecewiseLinearOn}, \texttt{hkDensity\_not\_isPiecewiseLinearPartition} and
\texttt{limit\_not\_isPiecewiseLinearOn}. The file contains no \texttt{sorry} and no \texttt{native\_decide}.

\end{document}
